\documentclass{article}

\usepackage{float}
\restylefloat{table}

\usepackage{booktabs}

\title{Team Productivity: POC\\\progname}

\author{\authname}

\date{}

\input{../Comments.text}
\input{../Common.text}

\begin{document}

\maketitle

This document summarizes the contributions of each team member up to the POC
Demo.  The time period of interest is the time between the beginning of the term
and the POC demo.

\noindent \textbf{Project URL: }
\url{https://github.com/JackEaston4/Troitsky-Toolbox}

\noindent Data in this report was collected from the repository's \texttt{main}
branch, issue tracker and pull requests on November 19, 2026, covering the
period from the start of term (September 8, 2026) to that date.

\section{Demo Plans}

\subsection*{Feasibility Risk}

The feasibility risk is whether the Toolbox can read a bridge's 3D CAD model
(a STEP file), display it in the browser, and extract usable measurements from
it.

This is the most complex single feature in the Toolbox, and the planned rule
compliance and bridge performance tools depend on the data it extracts. If it
is not achievable, those tools need a different source of bridge data.

Connecting the Toolbox's many independent tools is not treated as a POC risk,
since most tools are simple on their own; bridge visualization and data
extraction is the individually most complex part, so it is the focus of the
POC.

\subsection*{What Will Be Demonstrated}

A live demonstration of running code, showing that a bridge model can go from a
CAD file to a viewable, measured model in the Toolbox:
\begin{itemize}
\item A member uploads a bridge model exported from CAD.
\item The Toolbox checks that the file is a readable model, and rejects
  unreadable files with a message saying why.
\item The model is displayed in 3D, and the member can rotate and zoom it to
  confirm it was read correctly.
\item The Toolbox shows basic measurements taken from the model, such as its
  dimensions and estimated mass, and how long the model took to load.
\end{itemize}

The POC is not a minimum viable product. It shows that this pipeline works, so
the analysis tools that depend on it can be built on top of it.

If the demo shows that model extraction is not achievable, the team will change
the scope to an alternate feature that avoids the technical limitation, as
described in the Development Plan.

\section{Team Meeting Attendance}

\begin{table}[H]
\centering
\begin{tabular}{ll}
\toprule
\textbf{Student} & \textbf{Meetings}\\
\midrule
Total & 15\\
Pedro Azevedo & 15\\
Mason Azzopardi & 15\\
Michael Roberts & 13\\
Jack Easton & 15\\
\bottomrule
\end{tabular}
\end{table}

Each team meeting is recorded as a GitHub issue with the \texttt{meeting}
label, an attendance checklist and an agenda.

Counts are taken from the attendance checklists of the 15 \texttt{Team Meeting}
issues (\#5, \#30, \#51, \#87, \#96, \#120, \#149, \#182, \#221, \#224, \#237, \#284,
\#287, \#292, \#310).

At the meeting of October 28 (\#221), one member joined late and online; this
is counted as attended.

\section{Supervisor/Stakeholder Meeting Attendance}

\noindent \textbf{Supervisor's Name: } Isabella Vescio (Troitsky Club
Co-President)

\begin{table}[H]
\centering
\begin{tabular}{ll}
\toprule
\textbf{Student} & \textbf{Meetings}\\
\midrule
Total & 3\\
Pedro Azevedo & 3\\
Mason Azzopardi & 3\\
Michael Roberts & 0\\
Jack Easton & 0\\
\bottomrule
\end{tabular}
\end{table}

The repository has no supervisor meeting issues for this period. The counts
above are the three stakeholder interview sessions held with Troitsky Club
members for the Fall Custom Document (November 12, 13 and 14, 2026), as recorded
in the sessions table of the interview report. Each session was run by one
interviewer and one notetaker.

\section{Lecture Attendance}

\begin{table}[H]
\centering
\begin{tabular}{ll}
\toprule
\textbf{Student} & \textbf{Lectures}\\
\midrule
Total & 11\\
Pedro Azevedo & 1\\
Mason Azzopardi & 1\\
Michael Roberts & 5\\
Jack Easton & 5\\
\bottomrule
\end{tabular}
\end{table}

The total is the number of lectures on the course calendar between the start
of class and the POC demos (September 9, 10, 16, 17, 21, 23, 24 and 28, October
5 and 21, and November 2), excluding TRON sessions and tutorials. Lecture issues
exist for 5 of these lectures (September 24 and 28, October 5 and 21, November
2; issues \#31, \#63, \#95, \#207, \#229), and the counts above are taken from
their attendance checklists. Issue \#64 duplicates the September 28 lecture and
is not counted separately.

\section{TA Document Discussion Attendance}

\noindent \textbf{TA's Name: } Mohamad Mirhoseininejad (initial TA), Melvin
Mokhtari (from October 30, 2026)

\begin{table}[H]
\centering
\begin{tabular}{ll}
\toprule
\textbf{Student} & \textbf{Meetings}\\
\midrule
Total & 4\\
Pedro Azevedo & 3\\
Mason Azzopardi & 4\\
Michael Roberts & 4\\
Jack Easton & 4\\
\bottomrule
\end{tabular}
\end{table}

Counts are taken from the four \texttt{TA Meeting} issues (\#10, September 23;
\#99, October 8; \#223, October 30; \#286, November 11). The October 30 meeting
was held outside the allocated TA meeting time because of the TA's schedule.

\section{Commits, Lines Added and Lines Deleted, Pull Requests}

\begin{table}[H]
\centering
\setlength{\tabcolsep}{4pt}
\begin{tabular}{llllll}
\toprule
\textbf{Student} & \textbf{Commits} & \textbf{Percent} &
\begin{tabular}[b]{@{}l@{}}\textbf{Lines}\\\textbf{Added}\end{tabular} &
\begin{tabular}[b]{@{}l@{}}\textbf{Lines}\\\textbf{Deleted}\end{tabular} &
\begin{tabular}[b]{@{}l@{}}\textbf{Pull}\\\textbf{Requests}\end{tabular}\\
\midrule
Total & 340 & 100\% & 16745 & 12061 & 106\\
Pedro Azevedo & 125 & 36.8\% & 5519 & 7073 & 38\\
Mason Azzopardi & 101 & 29.7\% & 3586 & 1439 & 36\\
Michael Roberts & 49 & 14.4\% & 4511 & 2077 & 17\\
Jack Easton & 65 & 19.1\% & 3129 & 1472 & 15\\
\bottomrule
\end{tabular}
\end{table}

Commits are non-merge commits on \texttt{main}; the 139 merge commits
are excluded. A further 100 commits made by the
\texttt{github-actions[bot]} account, which regenerates the PDFs, are excluded.
Lines added and deleted exclude binary files and the initial commit that
imported the course template (11354 lines). Pull requests count every pull
request opened (open, merged and closed).

There is a significant discrepancy in lines of code committed into \texttt{main},
as well as lines added, lines deleted and pull requests. However, this should
not be regarded as evidence against the fact that each team member
has done their fair share. The discrepancy comes from Pedro Azzopardi
committing his vast interview records and other information specific to the expansive
knowledge he has gained as a veteran member of the Troitsky club. Mason has spearheaded AI-assisted
formatting which boosts the number of lines and pull requests, because virtually every doucment is heavily
formatted to conform to project standards. Jack created this repository, shaped many of the initial commit
counts, and worked extensively with Mason on the Module Guide. Michael's sudden rise in commits came from
his extensive travail on the Module Implementation Specification, including an initial draft based on
the component diagram which was later removed in favour of following the Module Guide.

\section{Issue Tracker}

\begin{table}[H]
\centering
\begin{tabular}{lll}
\toprule
\textbf{Student} & \textbf{Authored (O+C)} & \textbf{Assigned (C only)}\\
\midrule
Pedro Azevedo & 33 & 30 \\
Mason Azzopardi & 81 & 21 \\
Michael Roberts & 44 & 38 \\
Jack Easton & 32 & 40 \\
\bottomrule
\end{tabular}
\end{table}

A further 17 issues were authored by other teams' peer reviewers. Of the closed
issues, 30 have no assignee and are not counted in the Assigned column.

This metric is fair across the board. Mason created all the SRS issues for the group
while Michael created the milestones and an issue for each stage. However, the assignments
are equitable with only minor deviations. This can be explained by granularity, because sometimes
multiple related tasks are combined into one issue.

\section{CICD via GitHub Actions}

The repository uses one GitHub Actions workflow,
\href{https://github.com/JackEaston4/Troitsky-Toolbox/blob/main/.github/workflows/latex-pages.yml}{\texttt{latex-pages.yml}}.
On every push and pull request to \texttt{main} that changes the documentation,
it compiles the changed LaTeX documents and commits the generated PDFs back to
the repository for TA review.

Continuous Integration / Continious Delivery has been extensively used in this project and will continued to be
drawn upon in order to aid our continued success. All code is merged into our repository, namely the \texttt{main} branch,
after peer review to ensure quality. Automated tests ensure, and will ensure once implementation begins, that the 
code compiles and passes appropriate user tests, catching a whole swath of potential errors before being merged into
production. Manual review by an uninvolved developer (or extensive group-wide review when the entire group is involved)
augments our quality assurance proccess, to ensure all production changes improve the quality of production. Once a 
version of \texttt{main} is release-ready, all future versions of \texttt{main} must be release-ready as well.

To avoid duplicate work a developer shall assign themselves to the corresponding issue before commencing work on the task. 
If not such issue exists, the developer shall create one. 
The developers shall keep the Kanban board up to date and classify all issues by milestone. This will ensure 
issues are worked on in order of urgency and importance, instead of easiness, visibility, expected gratitude, or 
alignment with other arbitrary criteria.

\section{Custom Documents (CDs)}

The chosen CDs for the project are the Interview Report (Fall) and 
a Usability Report (Winter). Both will be kept moving forward, matching 
the project requirements. The CDs are linked in the project READMEs and 
contained in docs/CDs.

The CDs were split given the connection and required proximity and 
communication with the Troitsky club. Since Pedro Azevedo is a competing 
member and executive in Troitsky, he has served as the primarily lead for 
Fall CD (Interviews). Other members served to help analyze interview findings, 
help take notes during the interview and help write and review the report.

\section{Team Charter Trigger Items}

The only properly triggered item is listed below:
\begin{itemize}
\item Approaching a deliverable, a single missed scheduled meeting must be
  replaced with a supporting meeting or a direct discussion of the meeting
  contents.
\end{itemize}

There was no major infractions that would require next steps, the triggered item 
was for the sake of all members benefit.

\section{Additional Productivity Metrics}

No other additional project metrics. Since the project also entails integration with the McMaster Troitsky Club, 
there is work done that is not as clearly documented, but will be present and valuable to the project's success.


\end{document}
